% ═══════════════════════════════════════════════════════════════
% Section 1: Introduction
% ═══════════════════════════════════════════════════════════════

\section{Introduction}
\label{sec:introduction}

Brain-computer interfaces (BCIs) are no longer confined to research laboratories.
Neuralink has implanted its N1 device in human patients~\cite{musk2019neuralink},
Synchron's Stentrode has demonstrated motor neuroprosthesis via neurointerventional
surgery~\cite{oxley2021stentrode}, and Blackrock Neurotech's Utah array has enabled
high-performance speech neuroprostheses~\cite{willett2023speech}. Paradromics is
advancing toward high-bandwidth cortical interfaces for paralyzed patients.
Investment in neurotechnology grew 700\% between 2014 and 2021, with the global
market projected to reach \$25 billion by 2030~\cite{unesco2025recommendation}.

Regulation has begun to require cybersecurity for these devices without saying how to
assess it. Section 3305 of the Food and Drug Omnibus Reform Act of 2022
(FDORA, Pub.~L. 117-328) added Section 524B to the FD\&C Act~\cite{fda524b},
mandating cybersecurity documentation---including threat modeling, software bills
of materials, and vulnerability disclosure---in premarket submissions for
connected medical devices, and since October 2023 the FDA has applied a
Refuse-to-Accept policy to non-compliant submissions. Section 524B does not specify
a threat taxonomy or scoring system, and we are unaware of provisions in the standards
it draws on (CVSS, IEC 62443, AAMI TIR57, ISO 14971) for the risks specific to devices
that read and write neural signals. We found no neural-specific security requirements in
the EU Medical Device Regulation~\cite{eumdr2017}.

\subsection{Why Neurosecurity}
\label{sec:why-neurosecurity}

Four properties of neural devices, each documented in the literature, separate their
security from that of other connected medical devices.

\paragraph{The device acts on a person.}
A stimulating device delivers current to tissue. The same mechanisms underlie established
therapies such as deep brain stimulation~\cite{lozano2019dbs} and transcranial magnetic
stimulation~\cite{hallett2007tms}, so an unauthorised change to stimulation parameters is a
change to a treatment, not to a record. Some such effects may not be reversible, and a
security assessment has to say so.

\paragraph{Neural data cannot be reissued.}
A leaked password can be changed; a recording of a person's brain activity cannot. Present
devices do not read thoughts~\cite{ienca2018brainleaks,wexler2019neurohype}: speech
neuroprostheses decode attempted speech using models trained on a cooperating
user~\cite{willett2023speech}. Information can nevertheless leak. Martinovic et
al.~\cite{martinovic2012feasibility} showed that P300 responses recorded with a consumer
headset reveal information about PINs, banks and other private details at better than
chance. If decoders keep improving, a recording retained today becomes a readout risk later; the
catalogue records this as harvest-now-decode-later (QIF-T0166), with theoretical status.

\paragraph{The device can outlast its support.}
An implant stays in a person after the party that maintains it withdraws. In a modified
Delphi study with a 24-expert panel, the 14 experts who completed the final survey agreed
that responsibility for implanted neural devices after a trial needs explicit
rules~\cite{higgins2026posttrial}. No attacker is needed for harm
here, which is why this revision keeps such conditions in a separate hazards dataset
(Section~\ref{sec:tara}).

\paragraph{The attack surface spans biology and silicon.}
A recent review treats every element of the loop between brain, device, linked devices and
servers, and every link between them, as an attack surface, and describes defence frameworks
built specifically for BCIs as research prototypes~\cite{bagley2026threatvectors}.

Two consequences follow. First, an assessment needs dimensions that information-technology
scoring lacks (Section~\ref{sec:niss}). Second, the evidence base is thin: of the
\LiveTechniques{} techniques in the live catalogue, 95 are theoretical, and none of the
catalogue's status values means that an attack has been carried out on a person with a
device. Neurosecurity at this stage is therefore anticipatory work, and its credibility
depends on separating what has been shown from what is supposed. The evidence rules added
in this revision (Section~\ref{sec:tara}) exist for that reason.

\subsection{The Scoring Gap}

The Common Vulnerability Scoring System (CVSS v4.0)~\cite{first2023cvss4} is the
industry standard for assessing vulnerability severity. CVSS captures exploitability
characteristics---attack vector, complexity, privileges required, user
interaction---and impact on the confidentiality, integrity and availability of
information systems. It has no metric for:

\begin{itemize}
  \item \textbf{Biological tissue damage}---seizure induction, tissue injury,
        involuntary motor activation
  \item \textbf{Unauthorised reading or disruption of neural function}---inference
        from neural data, perceptual or motor disruption
  \item \textbf{Consent boundaries}---the distinction between operating within
        consented parameters and covert manipulation
  \item \textbf{Damage reversibility}---IT assets can be restored from backup;
        neural tissue cannot
  \item \textbf{Neuroplastic consequences}---prolonged stimulation can cause
        lasting change to neural pathways
\end{itemize}

When a vulnerability in a BCI could induce seizures, expose information inferred from
neural signals, or cause lasting harm, a CVSS base score alone is insufficient. In the
\SnapshotTechniques{}-technique snapshot analysed here, CVSS cannot express the primary
impact of 59 techniques and misses part of it for 30 more (Section~\ref{sec:niss}).

\subsection{Contributions}

This paper presents a proposed, unvalidated framework with four parts:

\begin{enumerate}
  \item \textbf{An addressing model} (Section~\ref{sec:hourglass}): eleven codes in
        use---seven for neural structures (N1--N7), one for the electrode-tissue
        interface (I0) and three for device systems (S1--S3)---each with a membership
        rule. The codes are addresses with no order. Earlier versions of this paper
        presented them as an hourglass stack with a chokepoint; Section~\ref{sec:hourglass}
        says why that was withdrawn.

  \item \textbf{TARA} (Section~\ref{sec:tara}): a catalogue of
        \SnapshotTechniques{}~techniques across \SnapshotTactics{}~tactics in the analysed
        snapshot (\LiveTechniques{} in the live registry), each classified by evidence
        status, severity and dual-use therapeutic potential. TARA is an independent
        taxonomy inspired by the structure of MITRE ATT\&CK\textsuperscript{\textregistered}.

  \item \textbf{NISS} (Section~\ref{sec:niss}): the Neural Impact Scoring
        System, an impact-only CVSS v4.0 extension designed to conform with
        FIRST.org's extension mechanism (User Guide
        \S3.11)~\cite{first2023cvss4userguide}, adding six neural-specific
        metrics, with stated limits and three scoring states.

  \item \textbf{The Neural Impact Chain} (Section~\ref{sec:nic}): a threat-modeling
        method that associates techniques with diagnostic category references through
        addresses, structures and functions to DSM-5-TR and ICD-10-CM
        codes~\cite{apa2022dsm5tr}. The codes are category references, not
        diagnostic claims.
\end{enumerate}

Section~\ref{sec:governance} sets out why a security framework needs neuroethics and
how it enters the scoring; Section~\ref{sec:case-studies} presents case studies, and
Section~\ref{sec:limitations} states limitations and what this revision changed. The
framework, catalogue and scoring system are released under the Apache 2.0 license.
